\subsection*{Scalars, eigenvalues, and eigenspaces}


\begin{problem}[Three quadratic characteristic equations]\label{ex:08:two-eigenvalues}
Find all eigenvalues, including nonreal ones, of
\[
 A=\begin{pmatrix}3&5\\-2&-4\end{pmatrix},\quad
 B=\begin{pmatrix}3&4\\-5&-5\end{pmatrix},\quad
 C=\begin{pmatrix}3&5\\5&7\end{pmatrix}.
\]
\end{problem}
\begin{solution}
Their characteristic polynomials are
\[
 \begin{aligned}
 p_A(t)&=(t-3)(t+4)+10=(t-1)(t+2),\\
 p_B(t)&=(t-3)(t+5)+20=(t+1)^2+4,\\
 p_C(t)&=(t-3)(t-7)-25=(t-5)^2-29.
 \end{aligned}
\]
Thus the eigenvalues are \(1,-2\); \(-1+2i,-1-2i\); and
\(5+\sqrt{29},5-\sqrt{29}\), respectively.
\end{solution}

\begin{problem}[Three eigenvalues with different block structures]\label{ex:08:three-eigenvalues}
Find the eigenvalues of
\[
 A=\begin{pmatrix}7&0&0\\0&3&0\\0&0&3\end{pmatrix},\quad
 B=\begin{pmatrix}7&0&0\\0&3&5\\0&5&4\end{pmatrix},\quad
 C=\begin{pmatrix}7&0&0\\0&3&1\\0&0&3\end{pmatrix}.
\]
\end{problem}
\begin{solution}
For \(A,C\), triangularity gives \(7,3,3\).
For \(B\), expansion in the first column gives
\[
 p_B(t)=(t-7)\{(t-3)(t-4)-25\}=(t-7)(t^2-7t-13).
\]
Its eigenvalues are \(7,(7+\sqrt{101})/2,(7-\sqrt{101})/2\).
The common list for \(A,C\) alone does not determine eigenspace dimensions.
\end{solution}

\begin{problem}[A repeated root and its eigenspace]\label{ex:08:ma-three}
For \(A=\begin{pmatrix}1&3&0\\0&1&0\\2&1&5\end{pmatrix}\), find the
characteristic equation, eigenvalues, rank, and all eigenvectors.
Can independent eigenvectors be chosen? Are the eigenspaces orthogonal?
\end{problem}
\begin{solution}
Expansion in the second row gives \(p_A(t)=(t-1)^2(t-5)\).
Thus the eigenvalues are \(1,1,5\), and \(\det A=5\), so its rank is \(3\).
For \(\lambda=1\), the equations are \(3y=0\), \(2x+y+4z=0\);
hence \(E_1=\operatorname{span}\{(-2,0,1)^\top\}\).
For \(\lambda=5\), \(-4y=0\) and \(-4x+3y=0\), giving
\(E_5=\operatorname{span}\{(0,0,1)^\top\}\).
All eigenvectors are the nonzero vectors in these two lines.
One vector from each gives an independent pair, but there is no
three-vector eigenbasis. Their dot product is \(1\), so the lines
are not orthogonal.
\end{solution}

\begin{problem}[The general real two-by-two matrix]\label{ex:08:quadratic}
For real \(A=\begin{pmatrix}a&b\\c&d\end{pmatrix}\):
\begin{enumerate}
\item Find \(p_A\) and all its eigenvalues, including nonreal ones.
Show that their sum and product are \(\tr A\) and \(\det A\).
\item Show that if one eigenvalue is nonreal, both are nonreal and
one is the conjugate of the other.
\item If \(b=c\), prove that both eigenvalues are real.
\end{enumerate}
\end{problem}
\begin{solution}
Set \(\Delta=(a-d)^2+4bc\). Completing the square gives
\[
 p_A(t)=t^2-(a+d)t+ad-bc
       =\left(t-\frac{a+d}{2}\right)^2-\frac{\Delta}{4}.
\]
For \(\Delta\geqslant0\), the eigenvalues are
\((a+d\pm\sqrt\Delta)/2\), including a repeated root when \(\Delta=0\).
For \(\Delta<0\), they are
\[
 \lambda_\pm=\frac{a+d\pm i\sqrt{-\Delta}}2.
\]
In either case \(\lambda_++\lambda_-=a+d=\tr A\), and
\[
 \lambda_+\lambda_-=\frac{(a+d)^2-\Delta}{4}=ad-bc=\det A.
\]
The negative-discriminant case gives both assertions in(b) explicitly.
For \(b=c\), \(\Delta=(a-d)^2+4b^2\geqslant0\), proving(c).
\end{solution}

\begin{problem}[What does an eigenvalue determine?]\label{ex:08:uniqueness}
\begin{enumerate}
\item Must matrices with the same characteristic polynomial be equal?
\item Can a nonzero vector be an eigenvector for two distinct eigenvalues
of the same matrix?
\item Can distinct vectors be eigenvectors for the same eigenvalue?
\end{enumerate}
\end{problem}
\begin{solution}
(a) No: \(I_2\) and \(\begin{pmatrix}1&1\\0&1\end{pmatrix}\) both have
polynomial \((t-1)^2\), but differ.
(b) If \(Av=\lambda v=\mu v\), then \((\lambda-\mu)v=0\); since \(v\ne0\),
\(\lambda=\mu\).
(c) Yes: \(v\) and \(2v\) have the same eigenvalue whenever \(v\) is an
eigenvector. For \(I_2\), the independent vectors \(e_1,e_2\) also share
an eigenvalue.
\end{solution}

\begin{problem}[A shear and a scalar matrix]\label{ex:08:shear}
Find all eigenvalues and eigenvectors of
\[
 A=\begin{pmatrix}1&1\\0&1\end{pmatrix},\qquad
 B=\begin{pmatrix}2&0\\0&2\end{pmatrix}.
\]
\end{problem}
\begin{solution}
For \(A\), the only eigenvalue is \(1\), with algebraic multiplicity two.
The equation \((A-I)(x,y)^\top=0\) is \(y=0\); its eigenvectors are
\((x,0)^\top\), \(x\ne0\). For \(B\), the only eigenvalue is \(2\),
also of multiplicity two, and every nonzero vector is an eigenvector.
Only \(B\) has a two-vector eigenbasis.
\end{solution}

\begin{problem}[Scaling, powers, and inverse]\label{ex:08:polynomials}
Let \(Av=\lambda v\), \(v\ne0\).
\begin{enumerate}
\item Show that \(tv\), \(t\ne0\), has the same eigenvalue, and that
\(v\) is an eigenvector of \(tA\) with eigenvalue \(t\lambda\).
\item Show that \(v\) is an eigenvector of \(sI-tA\) and of \(A^k\),
\(k\geqslant2\), and give their eigenvalues.
\item If \(A\) is invertible, prove the corresponding assertion for \(A^{-1}\).
\end{enumerate}
\end{problem}
\begin{solution}
Linearity gives \(A(tv)=\lambda(tv)\) and \((tA)v=(t\lambda)v\).
Also \((sI-tA)v=(s-t\lambda)v\).
Inductively \(A^kv=\lambda^kv\).
If \(A\) is invertible, \(\lambda=0\) would force \(v=0\); thus
\(\lambda\ne0\), and multiplying \(Av=\lambda v\) by \(A^{-1}\) gives
\(A^{-1}v=\lambda^{-1}v\).
The statement about \(tA\) allows \(t=0\); every nonzero vector is then
an eigenvector of the zero matrix.
\end{solution}

\begin{problem}[Transpose and eigenspaces]\label{ex:08:transpose}
For a real square matrix \(A\), show that \(A\) and \(A^\top\) have the
same characteristic polynomial and eigenvalues with multiplicities,
including nonreal ones.
Must they have the same eigenvectors?
\end{problem}
\begin{solution}
Transpose invariance of the determinant gives
\(\det(tI-A^\top)=\det((tI-A)^\top)=\det(tI-A)\).
Thus all roots, real or nonreal, and their multiplicities agree.
The eigenvectors need not coincide. For
\(A=\begin{pmatrix}1&1\\0&1\end{pmatrix}\),
\(\ker(A-I)=\operatorname{span}e_1\), whereas
\(\ker(A^\top-I)=\operatorname{span}e_2\).
\end{solution}

\begin{problem}[Conjugate pairs and odd dimension]\label{ex:08:odd-real}
\begin{enumerate}
\item Prove that nonreal eigenvalues of a real matrix occur in conjugate
pairs, with conjugate eigenvectors.
\item Prove that every real matrix of odd order has a real eigenvalue.
\end{enumerate}
\end{problem}
\begin{solution}
For \(Az=\lambda z\), \(z\ne0\), conjugate every coordinate:
\[
 \sum_j a_{ij}\bar z_j
 =\overline{\sum_j a_{ij}z_j}
 =\overline{\lambda z_i}
 =\bar\lambda\,\bar z_i.
\]
Here \(a_{ij}\) is real. Thus \(A\bar z=\bar\lambda\,\bar z\), with
\(\bar z\ne0\); when \(\lambda\) is nonreal, \(\bar\lambda\ne\lambda\).
For(b), the characteristic polynomial is monic of odd degree.
For large \(M>0\), \(p_A(M)>0\) and \(p_A(-M)<0\).
The intermediate value theorem gives a real \(\lambda\) with
\(p_A(\lambda)=0\). Then \(A-\lambda I\) has a nonzero real kernel vector.
This proves the odd-order assertion without needing a general theorem
counting all polynomial roots.
\end{solution}

\begin{problem}[*Eigenvalues of a real skew-symmetric matrix]\label{ex:08:skew}
Let \(A\) be real and \(A^\top=-A\). Prove that every eigenvalue is
zero or purely imaginary. What follows when the order of \(A\) is odd?
\end{problem}
\begin{solution}
Let \(A(x+iy)=(a+bi)(x+iy)\), with real \(x,y,a,b\) and
\((x,y)\ne(0,0)\). Equating real and imaginary parts gives
\[
 Ax=ax-by,\qquad Ay=bx+ay.
\]
For every real \(u\), skew-symmetry gives
\(u^\top Au=u^\top A^\top u=-u^\top Au\), hence \(u^\top Au=0\).
Take the dot products of the two equations with \(x,y\), then add:
\[
 0=x^\top Ax+y^\top Ay
  =a(\norm x^2+\norm y^2)+b(-x^\top y+y^\top x)
  =a(\norm x^2+\norm y^2).
\]
The factor in parentheses is positive, so \(a=0\).
Thus \(\lambda=bi\), with \(b=0\) allowed; in particular, the only
possible real eigenvalue is zero.
In odd order there is a real eigenvalue by the preceding problem.
It must be zero, so \(A\) is singular.
\end{solution}

\begin{problem}[Eigenvalues do not count the rank]\label{ex:08:shift}
Let \(J\) be \(n\times n\), with \(1\) immediately above the diagonal
and zero elsewhere. Find its eigenvalues and rank.
\end{problem}
\begin{solution}
It is triangular with zero diagonal, so \(p_J(t)=t^n\): all eigenvalues
are zero. Its action is
\(J(x_1,\ldots,x_n)^\top=(x_2,\ldots,x_n,0)^\top\).
Thus its image is \(\operatorname{span}(e_1,\ldots,e_{n-1})\),
of dimension \(n-1\), including rank zero for \(n=1\).
Counting nonzero eigenvalues would give the wrong rank for \(n>1\).
\end{solution}

\begin{problem}[An invertible perturbation]\label{ex:08:perturbation}
If \(A\) is singular, show that some nonzero real \(\epsilon\) makes
\(A+\epsilon I\) invertible.
\end{problem}
\begin{solution}
The polynomial \(q(\epsilon)=\det(A+\epsilon I)\) is monic of degree \(n\),
so it has only finitely many real roots. Choose any nonzero real number
outside that finite set. Then \(q(\epsilon)\ne0\).
The choice may be made arbitrarily close to zero, because every punctured
interval contains infinitely many real numbers.
\end{solution}

\begin{problem}[Similarity transports eigenspaces]\label{ex:08:similarity}
If \(B=P^{-1}AP\), prove that \(A,B\) have identical eigenvalues with
multiplicities. Must their eigenvectors be identical?
\end{problem}
\begin{solution}
\[
 \det(tI-B)=\det\{P^{-1}(tI-A)P\}=\det(tI-A).
\]
The eigenspace relation is \(E_\lambda(B)=P^{-1}E_\lambda(A)\), so
the actual vectors need not be the same.
For example \(A=\operatorname{diag}(1,2)\) and
\(P=\begin{pmatrix}1&1\\0&1\end{pmatrix}\) give
\(B=\begin{pmatrix}1&-1\\0&2\end{pmatrix}\).
The eigenvalue \(2\) has eigenvectors on the line \(\operatorname{span}e_2\)
for \(A\), but on \(\operatorname{span}(-1,1)^\top\) for \(B\).
\end{solution}

\begin{problem}[Rectangular products]\label{ex:08:ab-ba}
Let \(A\in\R^{m\times n}\) and \(B\in\R^{n\times m}\).
\begin{enumerate}
\item Show \(\det(I_m-AB)=\det(I_n-BA)\).
\item Show their nonzero eigenvalues, real or nonreal, agree with algebraic
multiplicities.
\item If \(n>m\), prove that \(BA\) has at least \(n-m\) zero eigenvalues.
Why may there be more?
\end{enumerate}
\end{problem}
\begin{solution}
Eliminate either off-diagonal block in
\(\begin{pmatrix}I_m&A\\B&I_n\end{pmatrix}\) to prove (a).
For \(t\ne0\), apply (a) with \(A/t\):
\[
 t^n p_{AB}(t)=t^m p_{BA}(t).
\]
This is a polynomial identity, since it holds for infinitely many \(t\).
Away from zero the powers of \(t\) do not affect root multiplicities,
giving (b). If \(n>m\), cancellation gives
\(p_{BA}(t)=t^{n-m}p_{AB}(t)\), proving (c).
There are additional zero roots when \(AB\) is singular; for \(A=0\),
all \(n\) eigenvalues of \(BA\) vanish.
\end{solution}

\begin{problem}[Elementary symmetric functions and power sums]\label{ex:08:newton-three}
Let \(\lambda_1,\lambda_2,\lambda_3\) be real and write
\[
 \prod_{i=1}^3(t-\lambda_i)=t^3+p_2t^2+p_1t+p_0,
\]
and put \(s_1=\sum_i\lambda_i\),
\(s_2=\lambda_1\lambda_2+\lambda_1\lambda_3+\lambda_2\lambda_3\),
\(s_3=\lambda_1\lambda_2\lambda_3\), and
\(\sigma_k=\sum_i\lambda_i^k\).
Find the \(p_i\) in terms of the \(s_i\), and prove that
\((s_1,s_2,s_3)\) and \((\sigma_1,\sigma_2,\sigma_3)\) determine one another.
\end{problem}
\begin{solution}
Expansion gives \(p_2=-s_1,p_1=s_2,p_0=-s_3\).
Furthermore
\[
 \sigma_1=s_1,\quad \sigma_2=s_1^2-2s_2,\quad
 \sigma_3=s_1^3-3s_1s_2+3s_3.
\]
The first follows by definition, the second by squaring \(s_1\),
and the third by expanding \(s_1^3\) and \(s_1s_2\).
Conversely
\[
 s_1=\sigma_1,\quad s_2=(\sigma_1^2-\sigma_2)/2,\quad
 s_3=(\sigma_3-\sigma_1^3+3\sigma_1s_2)/3.
\]
These explicit inverse formulas prove the one-to-one correspondence.
\end{solution}

\begin{problem}[When traces of powers determine eigenvalues]\label{ex:08:newton-general}
For real square matrices \(A,B\) of order \(n\), assume that both
characteristic polynomials split over \(\R\). Prove that their eigenvalues
agree with algebraic multiplicities if and only if
\(\tr A^k=\tr B^k\) for \(k=1,\ldots,n\).
\end{problem}
\begin{solution}
Triangularize a matrix as in the reading supplement. Its \(k\)-th power
has diagonal entries \(\lambda_i^k\), and trace is unchanged by
similarity, since \(\tr(XY)=\tr(YX)\). Thus
\(\tr A^k=\sum_i\lambda_i^k\).
To see that the first \(n\) power sums determine the characteristic
polynomial, write
\(f(z)=\prod_i(1-\lambda_i z)=\sum_{r=0}^n(-1)^r e_rz^r\), \(e_0=1\).
The product rule gives
\(f'(z)=-\sum_i\lambda_i\prod_{j\ne i}(1-\lambda_jz)\).
Multiply the finite geometric identity
\((1-\lambda_i z)\sum_{k=1}^n\lambda_i^kz^{k-1}
=\lambda_i-\lambda_i^{n+1}z^n\)
by \(\prod_{j\ne i}(1-\lambda_jz)\), and compare coefficients of
degrees below \(n\). Summing in \(i\) yields the recursion
\[
 r e_r=\sum_{k=1}^r(-1)^{k-1}e_{r-k}\sigma_k
 \quad(1\leqslant r\leqslant n).
\]
Hence the \(\sigma_k\) determine \(e_1,\ldots,e_n\) successively.
Equal traces give equal characteristic polynomials and therefore equal
root multisets. The converse follows directly from the power-sum formula.
\end{solution}

\begin{problem}[A matrix that cannot be diagonalized]\label{ex:08:three-jordan}
Prove that \(\begin{pmatrix}1&1&0\\0&1&1\\0&0&1\end{pmatrix}\)
cannot be diagonalized.
\end{problem}
\begin{solution}
Its characteristic polynomial is \((t-1)^3\).
The equation \((A-I)x=0\) gives \(x_2=x_3=0\), so its eigenspace
has dimension one and cannot provide a basis of \(\R^3\).
Equivalently, a diagonal matrix with all eigenvalues \(1\) is \(I\);
similarity to \(I\) would make \(A=PIP^{-1}=I\), a contradiction.
\end{solution}

\begin{problem}[Real QR factorization]\label{ex:08:qr}
Let \(A\in\R^{m\times n}\) have independent columns, \(m\geqslant n\).
Prove \(A=QR\), where \(Q^\top Q=I_n\) and \(R\) is real upper
triangular with positive diagonal.
\end{problem}
\begin{solution}
Apply the real Gram--Schmidt process to the columns \(a_j\).
Given \(q_1,\ldots,q_{j-1}\), put
\[
 r_{ij}=q_i^\top a_j\quad(i<j),\qquad
 w_j=a_j-\sum_{i<j}q_ir_{ij}.
\]
Taking the dot product with any earlier \(q_h\) gives
\(q_h^\top w_j=q_h^\top a_j-r_{hj}=0\).
The vector \(w_j\) is nonzero: otherwise \(a_j\) would lie in the span
of the preceding columns, contrary to independence.
Set \(r_{jj}=\norm{w_j}>0\), \(q_j=w_j/r_{jj}\), and \(r_{ij}=0\) for
\(i>j\). Induction gives orthonormal columns, and the identities
\(a_j=\sum_{i\leqslant j}q_ir_{ij}\) give \(A=QR\).
All quantities are real by construction.
\end{solution}

\begin{problem}[A rank-one product and its update]\label{ex:08:rank-one}
For real \(x,y\in\R^n\):
\begin{enumerate}
\item Prove that \(xy^\top\) has \(n-1\) zero eigenvalues and one
eigenvalue \(y^\top x\), counted with multiplicity.
\item Prove \(\det(I+xy^\top)=1+y^\top x\).
\item If \(\mu=1+y^\top x\ne0\), prove
\((I+xy^\top)^{-1}=I-\mu^{-1}xy^\top\).
\end{enumerate}
\end{problem}
\begin{solution}
The rectangular determinant identity, applied to an \(n\)-by-\(1\)
column and a \(1\)-by-\(n\) row, gives
\(p_{xy^\top}(t)=t^{n-1}(t-y^\top x)\), first for \(t\ne0\) and then
as a polynomial identity. This proves (a), including \(x=0\), \(y=0\),
or \(y^\top x=0\).
The same identity at \(I+xy^\top\) proves (b).
Put \(M=xy^\top\); then \(M^2=(y^\top x)M=(\mu-1)M\).
Both products of \(I+M\) and \(I-\mu^{-1}M\) equal
\(I+\{1-\mu^{-1}-\mu^{-1}(\mu-1)\}M=I\), proving (c).
\end{solution}

\begin{problem}[Left and right eigenvectors]\label{ex:08:left}
A real left eigenvector of \(A\) for a real \(\lambda\) is a nonzero
column \(y\) with \(y^\top A=\lambda y^\top\).
\begin{enumerate}
\item Show \(A^\top y=\lambda y\).
\item If \(Ax=\lambda x\), \(y^\top A=\mu y^\top\), and
\(\lambda\ne\mu\), prove \(y^\top x=0\).
\end{enumerate}
\end{problem}
\begin{solution}
Taking transposes proves (a). For (b), evaluate one scalar in two ways:
\[
 y^\top Ax=\lambda y^\top x=\mu y^\top x.
\]
Thus \((\lambda-\mu)y^\top x=0\), and the nonzero factor gives
\(y^\top x=0\).
This pairs a left eigenvector with a right eigenvector; two right
eigenvectors need not be orthogonal.
\end{solution}

\begin{problem}[A companion matrix]\label{ex:08:companion}
Let \(C\) have \(1\) in positions \((j,j+1)\), \(j<n\), last row
\((-a_0,-a_1,\ldots,-a_{n-1})\), and zero elsewhere.
Prove \(p_C(t)=t^n+a_{n-1}t^{n-1}+\cdots+a_1t+a_0\).
\end{problem}
\begin{solution}
For \(n=1\), \(C=[-a_0]\), giving \(t+a_0\).
For \(n>1\), expand \(tI-C\) in its first column.
The top entry contributes \(t\) times the determinant of the companion
pattern with coefficients \(a_1,\ldots,a_{n-1}\).
The bottom entry \(a_0\) contributes \(a_0\):
its cofactor sign is \((-1)^{n+1}\), and its minor is lower triangular
with \(n-1\) diagonal entries \(-1\), so the signs cancel.
Induction yields
\(t(t^{n-1}+a_{n-1}t^{n-2}+\cdots+a_1)+a_0\).
\end{solution}

\begin{problem}[Commuting matrices and simple eigenvalues]\label{ex:08:commuting-simple}
\begin{enumerate}
\item If \(AB=BA\) and \(\lambda\) is a simple real eigenvalue of \(A\),
prove that each eigenvector of \(A\) for \(\lambda\) is an eigenvector of \(B\).
\item If \(A\) has \(n\) distinct real eigenvalues and \(AB=BA\), prove that
one change of basis diagonalizes both matrices.
\end{enumerate}
\end{problem}
\begin{solution}
If \(Av=\lambda v\), then \(A(Bv)=B(Av)=\lambda Bv\).
The eigenspace \(E_\lambda(A)\) has dimension one because its geometric
multiplicity is positive and at most the algebraic multiplicity one.
Thus \(Bv=\mu v\) for a scalar \(\mu\), including \(\mu=0\).
Since \(v\ne0\), it is an eigenvector of \(B\).
In (b), the distinct eigenvalues of \(A\) provide a basis \(v_1,\ldots,v_n\).
By (a), the same vectors are eigenvectors of \(B\); their column matrix
diagonalizes both.
\end{solution}

\subsection*{Triangular form and polynomial consequences}
\begin{problem}[*Strict triangularity and rank decrease]\label{ex:08:nilpotent}
Let \(A\) be strictly upper triangular of order \(n\).
\begin{enumerate}
\item Prove \(A^n=0\).
\item Put \(r_j=\operatorname{rank}(A^j)\). Prove that
\(r_{j+1}<r_j\) whenever \(r_j>0\).
\end{enumerate}
\end{problem}
\begin{solution}
Every nonzero term in \((A^k)_{ij}\) requires a chain
\(i<i_1<\cdots<i_{k-1}<j\), so \(j-i\geqslant k\).
No such pair exists for \(k=n\), proving (a).
Let \(W=\operatorname{im}A^j\ne\{0\}\). Then \(AW=\operatorname{im}A^{j+1}\).
If the dimensions were equal, the map \(A:W\to W\) would be surjective
and hence invertible. Its \(n\)-th power would be invertible on the
nonzero space \(W\), contradicting \(A^n=0\).
Thus \(\dim AW<\dim W\), proving (b).
\end{solution}

\begin{problem}[A growing zero block]\label{ex:08:triangular-product}
Let \(A,B\) be upper triangular of order \(n\), \(1\leqslant k<n\).
Assume \(a_{ij}=0\) for \(1\leqslant i,j\leqslant k\), and
\(b_{k+1,k+1}=0\). Show that \(AB\) is upper triangular and its
upper-left \((k+1)\)-by-\((k+1)\) block is zero.
\end{problem}
\begin{solution}
An entry \(c_{ij}=\sum_\ell a_{i\ell}b_{\ell j}\) can have nonzero
terms only for \(i\leqslant\ell\leqslant j\); hence \(c_{ij}=0\)
when \(i>j\). If \(i\leqslant j\leqslant k\), every \(a_{i\ell}\)
in that range lies in the zero block. If \(j=k+1\), the terms with
\(\ell\leqslant k\) vanish from \(A\), and the remaining term contains
\(b_{k+1,k+1}=0\). This includes \(i=k+1\), proving the claim.
\end{solution}

\begin{problem}[An elementary similarity]\label{ex:08:elementary-similarity}
Let \(e_i,e_j\) be distinct standard basis vectors, and
\(T(\theta)=I-\theta e_ie_j^\top\).
\begin{enumerate}
\item Prove \(T(\theta)^{-1}=T(-\theta)\).
\item If \(A\) is upper triangular and \(i<j\), prove
\[
 T(\theta)^{-1}AT(\theta)=A+\theta(e_ie_j^\top A-Ae_ie_j^\top),
\]
and describe every possibly changed entry.
\item If \(a_{ii}\ne a_{jj}\), take
\(\theta=a_{ij}/(a_{ii}-a_{jj})\). What does this achieve?
\end{enumerate}
\end{problem}
\begin{solution}
Put \(E=e_ie_j^\top\). Since \(e_j^\top e_i=0\), \(E^2=0\);
thus \((I+\theta E)(I-\theta E)=I\) in either order.
Expansion in (b) gives \(A+\theta(EA-AE)-\theta^2EAE\).
Here \(EAE=a_{ji}E=0\), by upper triangularity.
For \(P=EA-AE\),
\[
 P_{h\ell}=\delta_{hi}a_{j\ell}-a_{hi}\delta_{j\ell}.
\]
The only possible changes are in row \(i\) at columns \(\ell\geqslant j\),
and in column \(j\) at rows \(h\leqslant i\).
At their intersection \(P_{ij}=a_{jj}-a_{ii}\); hence the chosen
\(\theta\) makes the new \((i,j)\) entry zero.
All other changed entries lie farther above the diagonal, and no
diagonal entry changes.
\end{solution}

\begin{problem}[Separate unequal diagonal blocks]\label{ex:08:block-separation}
Let an upper triangular \(A\) have distinct diagonal values
\(\lambda_1,\ldots,\lambda_k\), grouped contiguously with multiplicities
\(n_1,\ldots,n_k\). Write its diagonal blocks as \(A_1,\ldots,A_k\).
Using elementary similarities, prove that \(A\) is similar to
\(\operatorname{diag}(A_1,\ldots,A_k)\).
\end{problem}
\begin{solution}
An entry \((i,j)\) with \(i<j\) lying in two different diagonal blocks
has \(a_{ii}\ne a_{jj}\). Apply the preceding elementary similarity to
make it zero. That operation changes only positions \((i,\ell)\) with
\(\ell>j\), or \((h,j)\) with \(h<i\), besides the targeted entry.
These positions have strictly larger distance above the diagonal.
They also remain between different diagonal blocks, because
\(h\leqslant i<j\leqslant\ell\) and the blocks are contiguous.
Process distances \(j-i=1,2,\ldots,n-1\), removing all off-block
entries at each distance. No operation changes an entry already removed,
or any entry within a diagonal block. The resulting matrix is the
stated block diagonal matrix; the product of the elementary similarity
matrices is invertible.
\end{solution}

\begin{problem}[A unit triangular eigenvector matrix]\label{ex:08:unit-triangular}
If \(A\) is upper triangular with distinct diagonal entries, prove that
there is a unit upper triangular \(P\) such that
\(P^{-1}AP=\operatorname{diag}(a_{11},\ldots,a_{nn})\).
\end{problem}
\begin{solution}
For column \(j\), seek an eigenvector \(v_j\) for \(a_{jj}\) with
coordinates \(v_{jj}=1\) and \(v_{ij}=0\) for \(i>j\).
Rows below \(j\) and row \(j\) of \((A-a_{jj}I)v_j=0\) then hold.
For \(i=j-1,\ldots,1\), solve successively
\[
 (a_{ii}-a_{jj})v_{ij}=-\sum_{\ell=i+1}^{j}a_{i\ell}v_{\ell j}.
\]
The divisor is nonzero by distinctness. Thus all \(v_j\) exist.
Their column matrix \(P\) is upper triangular with diagonal \(1\), hence
invertible, and \(AP=P\operatorname{diag}(a_{ii})\).
\end{solution}

\begin{problem}[Triangular inverses]\label{ex:08:triangular-inverse}
\begin{enumerate}
\item Prove that an upper triangular matrix with nonzero diagonal has an
upper triangular inverse with nonzero diagonal.
\item Prove that a unit upper triangular matrix with integer entries has
a unit upper triangular inverse with integer entries.
\end{enumerate}
\end{problem}
\begin{solution}
In \(AX=I\), solve each column by back substitution.
For column \(j\), rows \(i>j\) give \(x_{ij}=0\), and row \(j\) gives
\(x_{jj}=a_{jj}^{-1}\). The remaining entries follow successively from
\[
 x_{ij}=-a_{ii}^{-1}\sum_{\ell=i+1}^j a_{i\ell}x_{\ell j}
 \quad(i<j).
\]
This proves (a). In (b), every divisor \(a_{ii}\) is \(1\);
the recurrence uses only sums and products of integers, proving (b).
Equivalently, if \(A=I+N\), strict triangularity gives \(N^n=0\), and
\(A^{-1}=I-N+\cdots+(-N)^{n-1}\).
\end{solution}

\begin{problem}[Trace of powers with real roots]\label{ex:08:trace-powers}
Let \(A\) be real and suppose
\(p_A(t)=\prod_{j=1}^n(t-\lambda_j)\), with real \(\lambda_j\).
Prove
\(\tr A^k=\sum_j\lambda_j^k\) for every integer \(k\geqslant1\).
\end{problem}
\begin{solution}
The real-root triangularization proof gives \(A=PTP^{-1}\),
where \(T\) is upper triangular with diagonal entries \(\lambda_j\).
Then \(A^k=PT^kP^{-1}\), and cyclic trace gives
\(\tr A^k=\tr(T^kP^{-1}P)=\tr T^k\).
For upper triangular \(U,V\), the only possible nonzero term in
\((UV)_{jj}=\sum_\ell u_{j\ell}v_{\ell j}\) has \(\ell=j\).
Thus \((UV)_{jj}=u_{jj}v_{jj}\). Induction gives
\((T^k)_{jj}=\lambda_j^k\); summing proves the formula.
\end{solution}

\begin{problem}[Nonzero eigenvalues and rank]\label{ex:08:rank-bound}
If the characteristic polynomial of a real \(n\times n\) matrix has
\(r\) nonzero roots, real or nonreal, counted with algebraic multiplicity,
prove that its rank is at least \(r\).
\end{problem}
\begin{solution}
Put \(d=\dim\ker A\). Extending a basis of \(\ker A\) to a basis of
\(\R^n\) represents \(A\) by
\(\begin{pmatrix}0&B\\0&C\end{pmatrix}\), where the first diagonal
block has order \(d\). Its characteristic polynomial is
\(t^d\det(tI-C)\). The polynomial on the right after removing \(t^d\)
has degree \(n-d\), so it can have at most \(n-d\) nonzero roots
with multiplicity. Hence \(r\leqslant n-d=\operatorname{rank}A\).
Equality need not hold, as a nonzero strictly upper triangular matrix shows.
\end{solution}

\begin{problem}[A simple eigenvalue and the converse]\label{ex:08:simple-rank}
\begin{enumerate}
\item If \(\lambda\) is a simple real eigenvalue, prove
\(\operatorname{rank}(\lambda I-A)=n-1\). Does the rank condition
imply that \(\lambda\) is simple?
\item Deduce the rank when zero is a simple eigenvalue.
If zero has algebraic multiplicity \(k\geqslant2\), must the rank be \(n-k\)?
\end{enumerate}
\end{problem}
\begin{solution}
A simple eigenvalue has a nonzero eigenspace of dimension at most one,
hence exactly one. Rank--nullity gives the stated rank.
Conversely the rank condition implies a nonzero kernel and thus an
eigenvalue, but it need not be simple:
\(\begin{pmatrix}0&1\\0&0\end{pmatrix}\) has rank one and a double
eigenvalue zero.
If zero is simple, the rank is \(n-1\).
For multiplicity \(k\geqslant2\), take
\(A=\operatorname{diag}(J_k,I_{n-k})\), where \(J_k\) is the shift matrix.
Then \(p_A(t)=t^k(t-1)^{n-k}\), but
\(\operatorname{rank}A=(k-1)+(n-k)=n-1>n-k\).
\end{solution}

\begin{problem}[*The characteristic equation evaluated at the matrix]\label{ex:08:cayley-hamilton}
For a real \(n\times n\) matrix, write
\(p_A(t)=t^n+c_{n-1}t^{n-1}+\cdots+c_0\).
Prove
\[
 p_A(A)=A^n+c_{n-1}A^{n-1}+\cdots+c_0I=0.
\]
If a factorization \(p_A(t)=\prod_j(t-\lambda_j)\) is given, allowing
real or nonreal roots,
deduce \(\prod_j(\lambda_j I-A)=0\).
\end{problem}
\begin{solution}
For \(n=1\) the identity is \(A-a_{11}I=0\).
For \(n>1\), the entries of \(\operatorname{adj}(tI-A)\) are
polynomials of degree at most \(n-1\). Write
\[
 \operatorname{adj}(tI-A)=\sum_{j=0}^{n-1}B_jt^j.
\]
The cofactor identity, valid also when \(tI-A\) is singular, gives
\[
 (tI-A)\sum_{j=0}^{n-1}B_jt^j=p_A(t)I.
\]
Comparing coefficients yields
\[
 B_{n-1}=I,\qquad B_{k-1}-AB_k=c_kI\ (1\leqslant k<n),
 \qquad -AB_0=c_0I.
\]
Multiply the middle identity on the left by \(A^k\), sum in \(k\),
and include the last identity and \(A^n=A^nB_{n-1}\). The result is
\[
 p_A(A)=A^nB_{n-1}
 +\sum_{k=1}^{n-1}(A^kB_{k-1}-A^{k+1}B_k)-AB_0=0,
\]
because every term cancels with its neighbor.
No factorization of \(p_A\) was needed.
When the displayed factorization exists, its factors are
polynomials in \(A\), so they commute and their product is \(p_A(A)\).
Changing the sign of all \(n\) factors proves the final assertion.
\end{solution}

\subsection*{Further complete calculations and applications}
\begin{problem}[A reflection]\label{ex:08:reflection}
Let \(H=\begin{pmatrix}\cos\theta&\sin\theta\\
\sin\theta&-\cos\theta\end{pmatrix}\).
Find a nonzero real \(v\) with \(Hv=v\), and prove that every vector
perpendicular to \(v\) satisfies \(Hw=-w\).
\end{problem}
\begin{solution}
Write \(c=\cos\theta,s=\sin\theta\).
If \(c\ne1\), take \(v=(s,1-c)^\top\), which is nonzero.
Multiplication gives
\(Hv=(cs+s-sc,s^2-c+c^2)^\top=(s,1-c)^\top\).
A perpendicular vector is a multiple of \(w=(-(1-c),s)^\top\);
\[
 Hw=(-c+c^2+s^2,-s+sc-cs)^\top=(1-c,-s)^\top=-w.
\]
If \(c=1\), then \(s=0\), \(H=\operatorname{diag}(1,-1)\);
take \(v=e_1\), and every perpendicular vector is a multiple of \(e_2\).
The two eigendirections describe reflection in the line spanned by \(v\).
\end{solution}

\begin{problem}[A commuting map preserves an eigenspace]\label{ex:08:commuting-space}
Let \(S,T:V\to V\) be linear and \(ST=TS\).
If \(v\) is an eigenvector of \(T\) for \(\lambda\), show that \(Sv\)
is also one for \(\lambda\), provided \(Sv\ne0\).
\end{problem}
\begin{solution}
\(T(Sv)=S(Tv)=S(\lambda v)=\lambda Sv\).
The nonzero condition is needed only for the word eigenvector.
Without it, the always valid assertion is \(S(E_\lambda(T))\subseteq
E_\lambda(T)\).
\end{solution}

\begin{problem}[Two two-by-two eigenspace calculations]\label{ex:08:lang-two}
Find the characteristic polynomials, eigenvalues and eigenspace bases of
\[
 A=\begin{pmatrix}1&2\\3&2\end{pmatrix},\qquad
 B=\begin{pmatrix}3&2\\-1&0\end{pmatrix}.
\]
\end{problem}
\begin{solution}
For \(A\), \(p_A=(t-1)(t-2)-6=(t-4)(t+1)\).
The first row gives bases \((2,3)^\top\) for \(4\) and
\((-1,1)^\top\) for \(-1\).
For \(B\), \(p_B=t(t-3)+2=(t-1)(t-2)\).
The second row gives \(x=-\lambda y\), so bases are
\((-1,1)^\top\) for \(1\) and \((-2,1)^\top\) for \(2\).
Each listed vector also satisfies the other row.
Every root is simple, so each nonzero vector spans its entire eigenspace.
\end{solution}

\begin{problem}[Three three-by-three eigenspace calculations]\label{ex:08:lang-three}
Find the characteristic polynomials, eigenvalues and eigenspace bases of
\[
 A=\begin{pmatrix}4&0&1\\-2&1&0\\-2&0&1\end{pmatrix},\quad
 B=\begin{pmatrix}1&-3&3\\3&-5&3\\6&-6&4\end{pmatrix},\quad
 C=\begin{pmatrix}3&1&1\\2&4&2\\1&1&3\end{pmatrix}.
\]
\end{problem}
\begin{solution}
For \(A\), expansion in column \(2\) gives
\((t-1)\{(t-4)(t-1)+2\}=(t-1)(t-2)(t-3)\).
Bases for \(E_1,E_2,E_3\) are respectively
\[
 (0,1,0)^\top,\qquad(-1,2,2)^\top,\qquad(-1,1,1)^\top.
\]
They satisfy \(Av=\lambda v\), and simplicity of the roots proves
completeness.
For \(B\), expansion gives \(p_B=(t+2)^2(t-4)\).
Here \(B+2I\) has all rows multiples of \((1,-1,1)\), so
\[
 E_{-2}=\{(x,y,z)^\top:x-y+z=0\}
 =\operatorname{span}\{(1,1,0)^\top,(-1,0,1)^\top\}.
\]
The vector \((1,1,2)^\top\) spans \(E_4\).
For \(C\), \(p_C=(t-2)^2(t-6)\), and
\(C-2I=\begin{pmatrix}1\\2\\1\end{pmatrix}(1,1,1)\).
Consequently
\[
 E_2=\operatorname{span}\{(-1,1,0)^\top,(-1,0,1)^\top\},\qquad
 E_6=\operatorname{span}\{(1,2,1)^\top\}.
\]
The two independent vectors for each double root exhaust its algebraic
multiplicity; the simple eigenspaces are one-dimensional. All three
matrices are diagonalizable despite the repeated roots in \(B,C\).
\end{solution}

\begin{problem}[Two one-dimensional eigenspaces]\label{ex:08:lang-repeated}
Find all eigenvalues and eigenspace bases of
\[
 A=\begin{pmatrix}2&-1\\1&0\end{pmatrix},\qquad
 B=\begin{pmatrix}2&0\\1&2\end{pmatrix}.
\]
Show that the eigenvectors of each matrix span a one-dimensional space.
\end{problem}
\begin{solution}
For \(A\), \(p_A=(t-1)^2\), and the equation
\((A-I)(x,y)^\top=0\) is \(x=y\).
Thus \(E_1=\operatorname{span}\{(1,1)^\top\}\).
For \(B\), \(p_B=(t-2)^2\), and
\((B-2I)(x,y)^\top=0\) is \(x=0\).
Thus \(E_2=\operatorname{span}\{(0,1)^\top\}\).
These are all eigenvalues and eigenvectors; both spans have dimension one.
\end{solution}

\begin{problem}[Real eigenspaces of a cyclic shift and a cubic]\label{ex:08:lang-cyclic}
Find the real eigenvalues and eigenspace bases of
\[
 A=\begin{pmatrix}0&1&0&0\\0&0&1&0\\0&0&0&1\\1&0&0&0\end{pmatrix},
 \qquad B=\begin{pmatrix}-1&0&1\\-1&3&0\\-4&13&-1\end{pmatrix}.
\]
\end{problem}
\begin{solution}
For \(A\), the eigenvector equations give
\(x_2=\lambda x_1,x_3=\lambda^2x_1,x_4=\lambda^3x_1\) and
\(x_1=\lambda^4x_1\). A nonzero solution needs \(x_1\ne0\).
The only real solutions of \(\lambda^4=1\) are \(1,-1\).
Their eigenspaces have respective bases
\((1,1,1,1)^\top\) and \((1,-1,1,-1)^\top\).

For \(B\), expansion gives
\[
 p_B(t)=t^3-t^2-t-2=(t-2)(t^2+t+1).
\]
Since \(t^2+t+1=(t+1/2)^2+3/4>0\), its only real eigenvalue is \(2\).
Rows two and one of \((B-2I)(x,y,z)^\top=0\) give \(x=y\), \(z=3y\).
The last row then vanishes, so the eigenspace is
\(\operatorname{span}\{(1,1,3)^\top\}\).
Neither matrix has enough real eigenvectors to form a basis.
\end{solution}

\begin{problem}[Trigonometric differentiation]\label{ex:08:derivative}
\begin{enumerate}
\item On the real space \(V=\operatorname{span}\{\sin t,\cos t\}\),
does differentiation have a nonzero eigenvector?
\item If \(k\) is a nonzero integer, show that \(\sin kx,\cos kx\)
are eigenvectors of the second derivative and give the eigenvalue.
\end{enumerate}
\end{problem}
\begin{solution}
In the ordered basis \((\sin t,\cos t)\), the derivative has matrix
\(\begin{pmatrix}0&-1\\1&0\end{pmatrix}\), with polynomial \(t^2+1\).
Thus it has no real eigenvalue and no nonzero eigenvector in the real
space \(V\).
Twice differentiating gives
\(D^2\sin kx=-k^2\sin kx\) and \(D^2\cos kx=-k^2\cos kx\).
Both functions are nonzero, and their common eigenvalue is \(-k^2\).
They may be considered in the invariant space they span, so no domain
ambiguity for the second derivative is needed.
\end{solution}

\begin{problem}[All eigenvectors in a distinct-eigenvalue basis]\label{ex:08:all-vectors}
Suppose \(v_1,\ldots,v_n\) is an eigenvector basis of \(T\), with
distinct eigenvalues \(c_1,\ldots,c_n\). Prove every eigenvector of \(T\)
is a scalar multiple of one of these basis vectors.
\end{problem}
\begin{solution}
Write \(v=\sum_i a_iv_i\) and suppose \(Tv=\lambda v\).
Independence of the basis gives \((c_i-\lambda)a_i=0\) for every \(i\).
At least one \(a_i\ne0\), so \(\lambda=c_i\) for that index.
Distinctness makes all other coefficients zero. Thus
\(v=a_iv_i\), with \(a_i\ne0\).
\end{solution}

\begin{problem}[Three further matrices]\label{ex:08:lang-further}
Find eigenvalues and eigenspace bases of
\[
 A=\begin{pmatrix}-1&2&2\\2&2&2\\-3&-6&-6\end{pmatrix},\quad
 B=\begin{pmatrix}3&2&1\\0&1&2\\0&1&-1\end{pmatrix},\quad
 C=\begin{pmatrix}-1&4&-2\\-3&4&0\\-3&1&3\end{pmatrix}.
\]
\end{problem}
\begin{solution}
Expansion gives
\[
 p_A=t(t+2)(t+3),\qquad p_B=(t-3)(t^2-3),\qquad
 p_C=(t-1)(t-2)(t-3).
\]
For \(A\), bases for eigenvalues \(-3,-2,0\) are respectively
\((-1,0,1)^\top,(-2,1,0)^\top,(0,-1,1)^\top\).
For \(B\), \(E_3=\operatorname{span}e_1\).
For \(\lambda=\pm\sqrt3\), the last row gives \(y=(\lambda+1)z\);
the first gives \(x=(2\lambda+3)z/(\lambda-3)\).
Taking \(z=1\) yields
\[
 E_\lambda=\operatorname{span}
 \left\{\left(\frac{2\lambda+3}{\lambda-3},\lambda+1,1\right)^\top\right\}.
\]
The second row holds because \(\lambda^2=3\).
For \(C\), bases for \(1,2,3\) are respectively
\((1,1,1)^\top\), \((2,3,3)^\top\), and \((1,3,4)^\top\).
For instance \(C(2,3,3)^\top=(4,6,6)^\top\).
All listed roots are simple; checking the displayed eigenvectors
therefore also proves that the bases exhaust the eigenspaces.
\end{solution}
